%
% 41-elementar.tex -- Elementarsymmetrische Polynome und der Fundamentalsatz
%
% (c) 2026 Prof Dr Andreas Müller
%
\subsection{Elementarsymmetrische Polynome und der Fundamentalsatz
\label{buch:komplex:symmetrisch:subsection:fundamentalsatz}}
Der Fundamentalsatz der Algebra besagt, dass sich ein Polynom in
Linearfaktoren zerlegen lässt.
Multipliziert man die Linearfaktoren aus, entstehen die Koeffizienten
des ursprünglichen Polynoms durch algebraische Ausdrücke aus den
Nullstellen.
In diesem Abschnitt soll untersucht werden, wie diese algebraischen
Ausdrücke entstehen.

%
% Permutationen
%
\subsubsection{Permutationen}
Im Folgenden soll untersucht werden, wie sich ein Polynom von mehreren
Variablen ändert, wenn die Variablen vertauscht werden.
Wir benötigen dazu die folgenden Notationen.

\begin{definition}[Permutation]
Eine \emph{Permutation} $\sigma$ von $n$ Elementen ist eine bijektive Abbildung
\index{Permutation}%
$\pi\colon [n]\to[n]$.
Die Menge aller Permutationen von $n$ Elementen bilden die
\emph{symmetrische Gruppe} $S_n$.
\index{symmetrische Gruppe}%
\index{Gruppe!symmetrisch}%
\index{Sn@$S_n$}%
\end{definition}

Sei jetzt $f(X_1,\dots,X_n)$ eine Funktion der $n$ Variablen 
$X_1,\dots,X_n$.
Die Vertauschung der Variablen mit einer Permutation $\pi\in S_n$
ergibt eine neue Funktion $f^\pi$, die durch
\[
f^\pi(X_1,X_2,\dots,X_n)
=
f(X_{\pi(1)},X_{\pi(2)},\dots,X_{\pi(n)})
\]
definiert ist.

\begin{beispiel}
Die Funktion $f(X_1,X_2) = \cos X_1\sin X_2 - \sin X_1\cos X_2$ ist eine
Funktion von zwei Variablen.
Die Gruppe $S_2$ enthält nur zwei Elemente, das neutrale Element und
die Vertauschung $\sigma\in S_2$ der beiden Variablen.
Die Funktion
\begin{align*}
f^\sigma(X_1,X_2)
&=
f(X_{\pi(1)},X_{\pi(2)})
=
f(X_2,X_1)
\\
&=
\cos X_2\sin X_1 - \sin X_2\cos X_1
\\
&=
-\cos X_1\sin X_2 + \sin X_1\cos X_2
=
-f(X_1,X_2),
\end{align*}
es ist also $f^\sigma=-f$.
Diese Aussage könnte man auch mit dem Additionstheorem bekommen, welches
\index{Additionstheorem}%
\begin{align*}
f(X_1,X_2)&=\sin(X_1-X_2)
\\
\Rightarrow\qquad
f^\sigma(X_1,X_2)
&=
\sin(X_2-X_1)
=
-\sin(X_1-X_2)
=
-f(X_1,X_2)
\end{align*}
garantiert.
\end{beispiel}

%
% Symmetrische Polynome
%
\subsubsection{Symmetrische Polynome}
Die Reihenfolge der Linearfaktoren, in die eine Polynom zerfällt, spielt
keine Rolle.
Jede Reihenfolge führt auf das gleiche ausmultiplizierte Polynom.
Um die Entstehung der Koeffizienten zu verstehen, müssen daher nur
algebraische Ausdrücke in den Nullstellen studiert werden, die sich nicht
ändern, wenn man die Nullstellen vertauscht.
Da beim Ausmultiplizieren nur Multiplikation und Addition verwendet
werden, müssen wir also Polynome in mehreren Variablen studieren, die
sich nicht ändern, wenn man die Variablen vertauscht.

\begin{definition}[symmetrisches Polynom]
\index{symmetrisches Polynom}%
Ein Polynom $p(X_1,\dots,X_n)$ in den Variablen $X_1,\dots,X_n$
heisst \emph{symmetrisch}, wenn jede Permutation $\sigma\in S_n$
der Variablen das gleiche Polynome ergibt.
Es gilt also
\[
p(X_{\sigma(1)},\dots,X_{\sigma(n)})
=
p(X_1,\dots,X_n)
\]
für alle Permutationen $\sigma\in S_n$.
\end{definition}

\begin{beispiel}
\label{buch:komplex:fundamental:symmetrisch:bsp:n=2}
Wir bestimmen die symmetrischen Polynome für $n=2$.
Es gibt nur zwei Permutationen für $n=2$, nämlich die identische
Permutation und die Vertauschung $1\leftrightarrow 2$.

Das Polynom $p(X_1,X_2)=X_1$ ist nicht symmetrisch, da die Vertauschung
$1\leftrightarrow 2$ daraus das Polynom $X_2\ne p(X_1,X_2)$ macht.
Erst die Kombination $\sigma_1(X_1,X_2)=X_1+X_2$ ist unter der Vertauschung
invariant.

Symmetrische Polynome zweiten Grades sind zum Beispiel
$\sigma_2(X_1,X_2)=X_1X_2$, aber auch
\begin{align*}
q(X_1,X_2)
&=
X_1^2+X_2^2.
\intertext{Allerdings ist dies kein ``neues'' Polynom, denn es ist}
X_1^2+X_2^2
&=
(X_1+X_2)^2 - 2X_1X_2
\\
&=
\sigma_1(X_1,X_2)^2 - 2 \sigma_2(X_1,X_2).
\end{align*}
Das Polynom $X_1^2+X_2^2$ kann also durch $\sigma_1(X_1,X_2)$ und
$\sigma_2(X_1,X_2)$ ausgedrückt werden.

Auch für symmetrische Polynome dritten Grades können wir Beispiele
finden.
In allen Fällen lassen sich diese Polynome durch die
$\sigma_k(X_1,X_2)$, $k=1,2$, ausdrücken:
\begin{align*}
X_1^2X_2 + X_1X_2^2
&=
(X_1+X_2)X_1X_2
=
\sigma_1(X_1,X_2)\cdot\sigma_2(X_1,X_2),
\\
X_1^3 + X_2^3
&=
(X_1+X_2)^3 - 3X_1^2X_2 - 3X_1X_2^2
\\
&=
(X_1+X_2)^3 - 3(X_1^2X_2 + X_1X_2^2)
\\
&=
\sigma_1(X_1,X_2)^3 - 3\sigma_1(X_1,X_2)\cdot \sigma_2(X_1,X_2).
\intertext{Dasselbe gilt für symmetrische Polynome vom Grad 4:}
X_1^2 X_2^2
&=
\sigma_2(X_1,X_2)^2
\\
X_1^3X_2+X_1X_2^3
&=
X_1X_2(X_1^2+X_2^2)
\\
&=
\sigma_2(X_1,X_2)\bigl(\sigma_1(X_1,X_2)^2-2\sigma_2(X_1,X_2)\bigr)
\\
&=
\sigma_1(X_1,X_2)^2 \sigma_2(X_1,X_2)
-2
\sigma_2(X_1,X_2)^2
\\
X_1^4+X_2^4
&=
\sigma_1(X_1,X_2)^4
-
4\sigma_1(X_1,X_2)^2 \sigma_2(X_1,X_2)
+8
\sigma_2(X_1,X_2)^2
-
6\sigma_2(X_1,X_2)^2
\\
&=
\sigma_1(X_1,X_2)^4
-
4\sigma_1(X_1,X_2)^2 \sigma_2(X_1,X_2)
+2
\sigma_2(X_1,X_2)^2.
\end{align*}
Keines der symmetrischen Polynome dritten oder vierten Grades enthält
Information, die sich nicht auch schon aus den beiden speziellen
Polynomen $\sigma_1(X_1,X_2)$ und $\sigma_2(X_1,X_2)$ ablesen lässt.
\end{beispiel}

Im Beispiel wurden zwei Polynome gefunden, mit denen sich alle 
anderen symmetrischen Polynome ausdrücken liessen.
Dies sind die sogenannten elemetarsymmetrischen Polynome,
die im nächsten Abschnitt eingeführt werden.

%
% Elementarsymmetrische Polynome
%
\subsubsection{Elementarsymmetrische Polynome}
Wir betrachten das Polynom
\[
p(T,X_1,\dots,X_n)
=
(T+X_1)\cdot\ldots\cdot (T+X_n)
\]
in den Variablen $T,X_1,\dots,X_n$.
Da das Produkt kommutativ ist, ist dieses Polynom symmetrisch in den
Variablen $X_1,\dots,X_n$.
Multipliziert man das Polynom aus und fasst die Terme mit gleicher
Potenz von $T$ zusammen erhält man
\begin{equation}
\renewcommand{\arraycolsep}{2pt}
\begin{array}{rclclclcl}
p(T,X_1,\dots,X_n)
&=&
T^n
&+&
\sigma_1(X_1,\dots,X_n)
T^{n-1}
&+&
\sigma_2(X_1,\dots,X_n)
T^{n-2}
&+&
\ldots
\\
& &
\ldots
&+&
\sigma_{n-1}(X_1,\dots,X_n) T
&+&
\sigma_n(X_1,\dots,X_n)
\\
&=&
\rlap{$\displaystyle \sum_{k=0}^n \sigma_k(X_1,\dots,X_n) T^{n-k} $}
\end{array}
\label{buch:komplex:symmetrisch:eqn:definition}
\end{equation}
mit Polynomen $\sigma_k(X_1,\dots,X_n)$, $k=0,\dots,n$.
Da $p(T,X_1,\dots,X_n)$ symmetrisch ist, gilt
\[
p(T,X_{\sigma(1)},\dots,X_{\sigma(n)})
=
p(T,X_1,\dots,X_n)
\] 
für jede Permutation $\sigma\in S_n$.
Durch Koeffizientenvergleich folgt, dass
\[
\sigma_k(X_{\sigma(1)},\dots,X_{\sigma(n)})
=
\sigma_k(X_1,\dots,X_n)
\]
für jeden Koeffizienten $\sigma_k$, $k=1,\dots,n$, und für jede
Permutation $\sigma\in S_n$.
Insbesondere sind die $\sigma_k(X_1,\dots,X_n)$ alle symmetrische
Polynome.

\begin{definition}[Elementarsymmetrische Polynome]
\label{buch:komplex:fundamental:def:elementarsymmetrisch}
Die durch \eqref{buch:komplex:symmetrisch:eqn:definition} definierten
Polynome $\sigma_k(X_1,\dots,X_n)\in \mathbb{Z}[X_1,\dots,X_n]$ 
heissen die \emph{elementarsymmetrischen Polynome}.
\index{elementarsymmetrische Polynome}%
\end{definition}

\begin{beispiel}
Für $n=2$ entstehen die elementarsymmetrischen Polynome aus dem
Produkt
\[
(T+X_1)(T+X_2)
=
T^2 +(X_1+X_2)T+X_1X_2
\quad\Rightarrow\quad
\left\{
\begin{aligned}
\sigma_1(X_1,X_2) &= X_1+X_2 \\
\sigma_2(X_1,X_2) &= X_1X_2.
\end{aligned}
\right.
\]
Dies sind genau die Polynome, die im
Beispiel~\ref{buch:komplex:fundamental:symmetrisch:bsp:n=2}
gefunden wurden.
\end{beispiel}

\begin{beispiel}
Für $n=3$ 
kann man aus dem Produkt
\[
(T+X_1)
(T+X_2)
(T+X_3)
=
T^3
+
(X_1+X_2+X_3)T^2
+(X_1X_2+X_1X_3+X_2X_3)T
+
X_1X_2X_3
\]
die elementarsymmetrischen Polynome
\begin{equation}
\begin{aligned}
\sigma_1(X_1,X_2,X_3) &= X_1+X_2+X_3\\
\sigma_2(X_1,X_2,X_3) &= X_1X_2 + X_1X_3 + X_2X_3\\
\sigma_3(X_1,X_2,X_3) &= X_1X_2X_3
\end{aligned}
\label{buch:komplex:fundamental:bsp:elementar3}
\end{equation}
ablesen.
\end{beispiel}

Die Beispiele zeigen auch, dass für beliebiges $n$ die
elementarsymmetrischen Polynome für $k=0$, $k=1$ und 
$k=n$
die besonders einfache Form
\begin{align*}
\sigma_0(X_1,\dots,X_n) &= 1 \\
\sigma_1(X_1,\dots,X_n) &= X_1+\ldots + X_n \\
\sigma_n(X_1,\dots,X_n) &= X_1 \cdots X_n
\end{align*}
haben.
Die anderen elementarsymmetrischen Polynome können als
\[
\sigma_k(X_1,\dots,X_n)
=
\sum_{1\le i_1<i_2<\dots<i_k\le n}
X_{i_1}
X_{i_2}
\cdots
X_{i_k}
\]
geschrieben werden.

Die elementarsymmetrischen Polynome können dazu verwendet werden, die
Koeffizienten eines Polynoms aus den Nullstellen zu berechnen.
Die Nullstellen eines solchen Polynoms sind algebraische Zahlen über
dem Koeffizientenkörper.
Oft sind alle Nullstellen ausserhalb des Koeffzientenkörpers.
Trotzdem sind die Werte der elementarsymmetrischen Funktionen, ausgewertet
auf den Nullstellen des Polynoms, wieder im Koeffizientenkörper.

\begin{beispiel}
Das Polynom $p(X) = X^2+pX+q$ hat die Nullstellen
\[
\alpha_{\pm}
=
-\frac{p}{2}
\pm
\!
\sqrt{ \frac{p^2}{4}-q }.
\]
Setzt man diese Nullstellen in die elementarsymmetrischen Funktionen
für $n=2$ ein, erhält man
\[
\renewcommand{\arraycolsep}{1.5pt}
\begin{array}{rccclcl}
\sigma_1(\alpha_+,\alpha_-)
&=&
\alpha_++\alpha_-
&=&
\displaystyle
-\frac{p}{2}
+
\!\sqrt{ \frac{p^2}{4}-q }
+
-\frac{p}{2}
-
\!\sqrt{ \frac{p^2}{4}-q }
&=&
-p,
\\[7pt]
\sigma_2(\alpha_+,\alpha_-)
&=&
\alpha_+\cdot \alpha_-
&=&
\displaystyle
\biggl(
-\frac{p}{2}
+
\!\sqrt{ \frac{p^2}{4}-q }
\biggr)
\biggl(
-\frac{p}{2}
-
\!\sqrt{ \frac{p^2}{4}-q }
\biggr)
\\[7pt]
& &
&=&
\displaystyle
\biggl(-\frac{p}{2}\biggr)^2
-
\biggl(
\!\sqrt{ \frac{p^2}{4}-q }
\biggr)^2
\\[7pt]
& &
&=&
\displaystyle
\frac{p^2}{4}
-
\frac{p^2}{4}+q
&=&
q.
\end{array}
\]
Die Werte der elementarsymmetrischen Funktionen von $\alpha_{\pm}$ sind
also wieder im Koeffizientenkörper.
\end{beispiel}

\begin{satz}
Sei $p(X) = p_0+p_1X+\dots +p_{n-1}X^{n-1} + X^n\in K[X]$ ein normiertes
mit Koeffizienten in einem Körper $K$ und seien $\alpha_1,\dots,\alpha_n$
alle Nullstellen, so dass 
\[
p(X)
=
(X-\alpha_1)(X-\alpha_2)\cdots(X-\alpha_n) =
\prod_{k=1}^n (X-\alpha_k).
\]
Dann ist 
\[
(-1)^k
\sigma_k(\alpha_1,\dots,\alpha_n) 
=
p_k
\]
für $k=1,\dots,n$ oder
\[
p(X)
=
\sum_{k=0}^n
(-1)^k
\sigma_k(\alpha_1,\dots,\alpha_n) X^{n-k}.
\]
Ausserdem ist
\[
\sigma_k(\alpha_1,\dots,\alpha_n) \in K
\]
für $k=1,\dots,n$.
\end{satz}

\begin{beispiel}
Für das Polynom $p(X) = X^5 + 2X^4+3X^3+4X^2+5X+6\in\mathbb{Z}[X]$
kann man numerisch (zum Beispiel mit Maxima) die Nullstellen
\[
\renewcommand{\arraycolsep}{1.5pt}
\begin{array}{rcrcl}
\alpha_1 &=&  0.5516854634589816 &+& 1.2533488602772058 i\\
\alpha_2 &=&  0.5516854634589816 &-& 1.2533488602772058 i\\
\alpha_3 &=& -0.8057864693890309 &+& 1.2229047133744104 i\\
\alpha_4 &=& -0.8057864693890309 &-& 1.2229047133744104 i\\
\alpha_5 &=& -1.4917979881399015 & &
\end{array}
\]
finden.
Die Werte der elementarsymmetrischen Polynome müssen wieder
ganze Zahlen ergeben. 
Das Polynom $\sigma_1(X_1,\dots,X_5)$ ist die Summe
und $\sigma_5(X_1,\dots,X_5)$ ist das Produkt der Argumente.
Tatsächlich sind Summe und Produkt
\[
\renewcommand{\arraycolsep}{1.5pt}
\begin{array}{rcccl}
\sigma_1(\alpha_1,\dots,\alpha_5)
&=&
\displaystyle
\sum_{k=1}^5 \alpha_k
&=&
-2
\\
\sigma_5(\alpha_1,\dots,\alpha_5)
&=&
\displaystyle
\prod_{k=1}^5 \alpha_k
&=&
-6.000000000000004 + 3.525489416408166\cdot 10^{-16}i
\end{array}
\]
mit Maschinengenauigkeit bis auf das Vorzeichen die Koeffizienten
des Polynoms $p(X)$.
\end{beispiel}

%
% Der Fundamentalsatz für symmetrische Polynome
%
\subsubsection{Der Fundamentalsatz für symmetrische Polynome}
Die elementarsymmetrischen Polynome sind in einem gewissen Sinne
die einzigen symmetrischen Polynome, die man braucht.
Dies ist der Inhalt des folgenden Fundamentalsatzes für
symmetrische Polynome: jedes symmetrische Polynom lässt sich
als Polynom in den elementarsymmetrischen Polynomen ausdrücken.

\begin{satz}[Fundamentalsatz für symmetrische Polynome]
\label{buch:komplex:fundamental:satz:fundamentalpolynom}
Ist $p(X_1,\dots,X_n)\in K[X_1,\dots,X_n]$ ein symmetrisches
Polynom von $n$ Variablen, dann gibt es genau ein Polynom
$q(S_1,\dots,S_n)\in K[S_1,\dots,S_n]$ derart, dass
\index{Fundamentalsatz für symmetrische Polynome}%
\[
p(X_1,\dots,X_n)
=
q\bigl(\sigma_1(X_1,\dots,X_n),\dots,\sigma_n(X_1,\dots,X_n)\bigr).
\qedhere
\]
\end{satz}

Nicht nur gibt es das Polynom $q$, wir werden einen Beweis
geben, der das Polynom sogar rekursiv konstruiert.
Wir illustrieren die Idee am Beispiel eines symmetrischen
Polynoms $q(X_1,X_2)$ dritten Grades in $X_1$ und $X_2$.
Es muss mindestens einen der Terme
\[
X_1^3,\;
X_1^2X_2,\;
X_1X_2^2,\;
X_2^3
\]
dritten Grades enthalten.
Falls der Term $cX_1^3$ in $q$ vorkommt, dann muss auch der Term $cX_2^3$
vorkommen.
Mit dem elementarsymmetrischen Polynom $\sigma_1(X_1,X_2)$ können
wir ebenfalls ein Polynom dritten Grades konstruieren, welches
$cX_1^3$ und $cX_2^3$ enthält, nämlich $c\sigma_1(X_1,X_2)^3$.
Die Differenz $q_1(X_1,X_2)=q(X_1,X_2)-c\sigma_1(X_1,X_2)^3$ enthält die
Terme $X_1^3$ und $X_2^3$ nicht mehr.

$q_1(X_1,X_2)$ kann aber immer noch $c_1X_1^2X_2$ enthalten.
Symmetrie verlangt dann, dass es auch $c_1X_1X_2^2$ enthält.
Mit den elementarsymmetrischen Polynomen können wir 
\[
c_1\sigma_1(X_1,X_2)\sigma_2(X_1,X_2)
=
c_1(X_1+X_2)(X_1X_2)
=
c_1(X_1^2X_2+X_1X_2^2)
\]
konstruieren, welches die gleichen Terme enthält.
Das neue Polynom
$q_2(X_1,X_2) = q_1(X_1,X_2) - c_1\sigma_1(X_1,X_2)\sigma_2(X_1,X_2)$
enthält dann auch diese Terme nicht mehr.

Da es keine weiteren Terme dritten Grades gibt, ist $q_2$ jetzt nur noch
ein Polynom zweiten Grades.
Falls es den Term $c_2X_1^2$ enthält, können wir diesen durch
$q_3(X_1,X_2)=q_2(X_1,X_2)-c_2\sigma_1(X_1,X_2)^2$ eliminieren.
Falls dann noch ein Term $c_3X_1X_2$ bleibt, wird dieser durch
$q_4(X_1,X_2)=q_3(X_1,X_2)-c_3\sigma_2(X_1,X_2)$ eliminiert.

$q_4(X_1,X_2)$ ist ein Polynom ersten Grades und kann daher sofort
als Linearkombination von $\sigma_1(X_1,X_2)=X_1+X_2$ und
$\sigma_0(X_1,X_2)=1$ dargestellt werden.

Der Algorithmus eliminiert also die verschiedenen Arten von Monomen, die
in $q(X_1,X_2)$ vorkommen können dadurch, dass ein geeignetes Produkt von 
elementarsymmetrischen Polynomen subtrahiert wird.
Dabei ist in einer Reihenfolge vorzugehen, die verhindert, dass bereits
eliminierte Terme wieder hinzugefügt werden.
Es muss also als vor allem eine geeignete Anordnung der Monome definiert
werden.

Um die Ordnung zweier Monome gleichen Grades zu definieren, kürzen wir 
das Monom $X_1^{a_1}\dots X_n^{a_n}$ durch das $n$-Tupel
\[
e(X_1^{a_1}\cdots X_n^{a_n})
=
(a_n,a_{n-1},\dots,a_2,a_1)
\]
ab.
Die Funktion $e$ extrahiert also die Exponenten der einzelnen 
Variablen und ordnet sie nach absteigendem Index der Variablen an.
Zum $n$-Tupel $(a_n,\dots,a_1)$ ist $m=X_1^{a_1}\cdots X_n^{a_n}$
das Monom mit $e(m)=(a_n,\dots,a_1)$.

\begin{definition}[umgekehrt lexikographische Ordnung]
\label{buch:komplex:symmetrisch:def:ulexiko}
\index{umgekehrt lexikographische Ordnung}%
Die umgekehrt lexikographische Ordnung von $n$-Tupeln ist definiert durch
\[
(a_n,\dots,a_1) < (b_n,\dots,b_1)
\quad
:\Leftrightarrow
\quad
\renewcommand{\arraycolsep}{2pt}
\left\{\begin{array}{rclcl}
a_n &>& b_n	&\quad&\text{falls $a_n\ne b_n$}\\
a_{n-1} &>& b_{n-1}&&\text{falls $a_{n-1}\ne b_{n-1}$ und $a_n=b_n$}\\[-2pt]
&\vdots&        &&\vdots\\
a_2 &>& b_2	&&\text{falls $a_2\ne b_2$ und $a_k=b_k$ für $k\ge 3$}\\
a_1 &>& b_1	&&\text{falls $a_1\ne b_1$ und $a_k=b_k$ für $k\ge 2$.}
\end{array}
\right.
\]
Zwei Monome $m_1$ und $m_2$ gleichen Grades erfüllen $m_1 < m_2$, wenn
$e(m_1) < e(m_2)$.
\end{definition}

Die umgekehrt lexikographische Ordnung der Tripel vom Grad 3 ist
\index{lexikographisch}%
\begin{equation}
\begin{aligned}
(3,0,0)
&<
(2,1,0)
<
(2,0,1)
<
(1,2,0)
<
(1,1,1)
\\
&<
(1,0,2)
<
(0,3,0)
<
(0,2,1)
<
(0,1,2)
<
(0,0,3).
\end{aligned}
\label{buch:komplex:symmetrisch:eqn:3tupelordnung}
\end{equation}
``Grosse'' Elemente im Sinn der umgekehrt lexikographischen Ordnung
haben also möglichst grosse Komponenten $a_i$ mit kleinem Index $i$.

\begin{lemma}
\label{buch:komplex:symmetrisch:lemma:maximal}
Ein $n$-Tupel $(a_n,\dots,a_1)$ ist genau dann das grösste im Sinne
der umgekehrt lexikographischen Ordnung, wenn
$a_n\le a_{n-1}\le \ldots \le a_1$
ist.
\end{lemma}

\begin{proof}
Wäre $a_{i+1} > a_i$, dann liesse sich ein grösseres $n$-Tupel durch
Vertauschen von $a_i$ und $a_{i+1}$ finden:
\[
(a_n,\dots,a_{i+1},a_i,\dots,a_1)
<
(a_n,\dots,a_i,a_{i+1},\dots,a_1).
\]
Dies zeigt, dass das $n$-Tupel $(a_n,\dots,a_1)$ nicht das grösste sein 
kann, wenn die $a_k$ nicht monoton sind.
\end{proof}

Ein $n$-Tupel wie in Lemma~\ref{buch:komplex:symmetrisch:lemma:maximal}
wird auch ein \emph{maximales $n$-Tupel} genannt.
\index{maximales n-Tupel@maximales $n$-Tupel}%

\begin{definition}[Monomordnung]
\index{Monomordnung}%
Für zwei beliebige Monome $m_1$ und $m_2$ von möglicherweise verschiedenem
Grad wird die Ordnung durch
\[
m_1 < m_2
\quad
:\Leftrightarrow
\quad
\begin{cases}
\deg m_1 < \deg m_2&\qquad\text{falls $\deg m_1\ne \deg m_2$}\\
m_1 < m_2&\qquad\text{falls $\deg m_1 = \deg m_2$}
\end{cases}
\]
definiert, wobei die Ordnung gleichgradiger Monome im zweiten Fall
die ungekehrt lexikographische Ordnung von
Definition~\ref{buch:komplex:symmetrisch:def:ulexiko} ist.
\end{definition}

Für zwei Variablen $X_1$ und $X_2$ ergeben sich bis zum Grad 3 die folgenden 
Reihenfolgen:
\begin{align}
&\text{Grad 1:} & &X_1 < X_2
\notag
\\
&\text{Grad 2:} & &X_1^2 < X_1X_2 < X_2^2
\notag
\\
&\text{Grad 3:} & &
X_3^3
<
X_3^2 X_2
<
X_3^2 X_1
<
X_3 X_2^2
<
X_3 X_2 X_1
<
X_3 X_1^2
<
X_2^3
<
X_2^2
X_1
<
X_2
X_1^2
<
X_1^3
\label{buch:komplex:symmetrisch:eqn:3monomordung}
\end{align}
wobei sich
\eqref{buch:komplex:symmetrisch:eqn:3monomordung}
aus
\eqref{buch:komplex:symmetrisch:eqn:3tupelordnung}
ergibt.

\begin{definition}[führendes Monom, führender Koeffizient]
Das \emph{führende Monom} oder \emph{leading monimial}
$\operatorname{LM}(p)$ eines Polynoms $p$ ist
\index{fuhrendes Monom@führendes Monom}%
\index{leading monomial}%
\index{LM}%
\index{Monom!fuhrend@Monom!führend}%
das grösste Monom, welches in $p$ vorkommt.
Der \emph{führende Koeffizient} oder \emph{leading coefficient}
$\operatorname{LC}(p)$ von $p$ 
ist der Koeffizient des führenden Monoms in $p$.
\index{fuhrendes Koeffizient@führendes Koeffizient}%
\index{leading coefficient}%
\index{LC}%
\index{Koeffizient!fuhrend@Koeffizient!führend}%
\end{definition}

\begin{lemma}
Für $k=1,\dots,n$ gilt
\[
\operatorname{LM}(\sigma_k(X_1,\dots,X_n))
=
X_1\cdots X_k
\quad\text{mit}\quad
e(\operatorname{LM}(\sigma_k(X_1,\dots,X_n)))
=
(\underbrace{0,\dots,0}_{n-k},\underbrace{1,\dots,1}_{k}).
\]
\end{lemma}

\begin{proof}
Die Monome eines elementarsymmetrischen Polynoms enthalten keine
höheren als erste Potenzen jeder Variable.
In der definierten Ordnung hat das führend Monom daher nur die $k$
``grössten'' Variablen.
\end{proof}

\begin{lemma}
Ist $(a_n,\dots,a_1)$ ein maximales $n$-Tupel, dann ist
\[
m
=
\operatorname{LM}(
\sigma_1^{a_1-a_2}
\sigma_2^{a_2-a_3}
\cdots
\sigma_{n-1}^{a_{n-1}-a_n}
\sigma_{n}^{a_n}
)
\]
ein Monom mit $e(m)=(a_n,\dots,a_1)$.
\end{lemma}

\begin{proof}
Statt mit den Monomen zu rechnen, können wir mit den $n$-Tupeln rechnen:
\begin{align*}
\left.
\renewcommand{\arraycolsep}{2pt}
\begin{array}{crcl}
 & (0,\dots,0,0,1) &\cdot & (a_1-a_2)     \\
+& (0,\dots,0,1,1) &\cdot & (a_2-a_3)     \\
+& (0,\dots,1,1,1) &\cdot & (a_3-a_4)     \\[-2pt]
 &		   &\vdots&               \\
+& (0,\dots,1,1,1) &\cdot & (a_{n-1}-a_n) \\
+& (1,\dots,1,1,1) &\cdot & \phantom{(}a_n\\
\end{array}
\right\}
&=
\left\{
\renewcommand{\arraycolsep}{1pt}
%\begin{array}{clccclclcll}
\begin{array}{ccccccccccl}
 &(0       &,& \dots &,& 0           &,&       0     &,& a_1-a_2     &)\\
 &(0       &,& \dots &,& 0           &,& a_2-a_3     &,& a_2-a_3     &)\\
 &(0       &,& \dots &,& a_3-a_2     &,& a_3-a_4     &,& a_3-a_4     &)\\[-2pt]
 &\;\vdots & & \vdots& & \vdots      & & \vdots      & & \vdots      & \\
 &(0       &,& \dots &,& a_{n-1}-a_n &,& a_{n-1}-a_n &,& a_{n-1}-a_n &)\\
 &(a_n     &,& \dots &,& a_n         &,&  a_n        &,& a_n         &)
\end{array}
\right.
\\
&=
(a_n,\dots,a_3,a_2,a_1),
\end{align*}
wie verlangt.
\end{proof}

Nach diesen Vorbereitungen sind wir jetzt in der Lage, den
Satz~\ref{buch:komplex:fundamental:satz:fundamentalpolynom}
dadurch zu beweisen, dass wir den in den Vorbemerkungen zum Beweis
skizzierten Algorithmus ausformulieren.

\begin{proof}[Beweis von Satz~\ref{buch:komplex:fundamental:satz:fundamentalpolynom}]
Das Polynom $q$ wird iterativ durch den folgenden Algorithmus bestimmt:
\begin{enumerate}
\item Setze $q=0$.
\item Setze $a=e(\operatorname{LM}(p))$.
\item Ersetze $p$ und $q$ durch
\begin{align*}
p&\leftarrow
p - \operatorname{LC}(p) 
\sigma_1^{a_2-a_1}\sigma_2^{a_3-a_1}\cdots\sigma_{n-1}^{a_n-a_{n-1}},
\\
q&\leftarrow
q+ \operatorname{LC}(p) 
\sigma_1^{a_2-a_1}\sigma_2^{a_3-a_1}\cdots\sigma_{n-1}^{a_n-a_{n-1}}.
\end{align*}
\item Falls $p\ne 0$ weiter bei 2.
\end{enumerate}
Im Schritt~3 wird $p$ durch ein Polynom mit kleinerem führendem 
Monom ersetzt.
Nach endlich vielen Schritten terminiert der Algorithmus daher und
$q$ ist das gesuchte Polynom in den elementarsymmetrischen Funktionen.
\end{proof}

\begin{beispiel}
Wir führen den Algorithmus im Fall $n=3$ für das Polynom
\[
p(X_1,X_2,X_3) = {\color{darkred}X_1^3} + X_2^3 + X_3^3
\]
durch.
\smallskip

\noindent
Das führende Monom von $p$ ist 
\(
\operatorname{LM}(p) = {\color{darkred}X_1^3}
\)
mit führendem Koeffizienten $\operatorname{LC}(p)=1$,
somit sind die neuen Polynome $p$ und $q$
\begin{align*}
p
&\leftarrow
p-\sigma_1^3
= 
- 6 X_3 X_2 X_1
- 3(
  X_3^2 X_2
+ X_3^2 X_1
+ X_3 X_2^2
+ X_3 X_1^2
+ X_2^2 X_1
+ {\color{darkred}X_2 X_1^2}
)
\\
q
&\leftarrow
\sigma_1^3.
\intertext{Für das neue Polynom ist das führende Monom das rot
hervorgehobene $\operatorname{LM}(p)=X_2X_1^2$, mit
führendem Koeffizienten $\operatorname{LC}(p)=-3$, es muss als im nächsten
Schritt $3\sigma_1\sigma_2$ addiert werden:}
p
&\leftarrow
p+3\sigma_1\sigma_2
=
3 {\color{darkred}X_3X_2X_1}
\\
q
&\leftarrow
\sigma_1^3-3\sigma_1\sigma_2.
\intertext{Das führende Monom von $p$ ist jetzt
$\operatorname{LM}(p)={\color{darkred}X_3X_2X_1}=\sigma_3$ mit
führendem Koeffizienten $\operatorname{LC}(p)=3$,
sodass im nächsten Schritt $3\sigma_3$
subtrahiert werden muss:}
p
&\leftarrow
p - 3\sigma_3 = 0
\\
q
&\leftarrow
\sigma_1^3-3\sigma_1\sigma_2+3\sigma_3.
\end{align*}
Damit haben wir
\begin{equation}
X_1^3+X_2^3+X_3^3
=
\sigma_1(X_1,X_2,X_3)^3
-
3\sigma_1(X_1,X_2,X_3)\sigma_2(X_1,X_2,X_3)
+
3\sigma_3(X_1,X_2,X_3)
\label{buch:komplex:symmetrisch:eqn:X33}
\end{equation}
gefunden.
\end{beispiel}
